\paragraph{Measured the way motif finding measures, the picture flips.} If we count only the
\emph{strong} activations --- the ones above the $90$th-percentile threshold that discovery acts
on, not every nonzero --- the filters are already fairly separated, and the operator helps. With
the operator off, occupancy is just $2.6$ (of $50$ filters), and a quarter of active positions are
already private to a single filter (exclusivity $0.25$). Turn the operator on and it improves:
occupancy drops to $2.2$ and exclusivity climbs to $0.41$ --- now two-fifths of positions belong
to one filter.

\paragraph{Why the earlier read was too harsh.} Counting \emph{every} nonzero activation, the
same off layer looked hopeless (occupancy $\sim22$, half the bank on every spot). But most of
those are weak activations that motif finding ignores. Once you look only at the strong ones, the
crowding is far milder and the operator's nudge is clearly in the right direction --- more private
positions, fewer shared. So on the metric that actually matters downstream, the plain operator
already earns its keep; Experiment~3 asks how much further the sharper modes can push it.
